\documentclass[11pt]{article}
\usepackage{amsmath,amssymb,geometry}
\geometry{margin=1in}
\title{ROOT V136: Geometric and Yukawa Null Tests (Model, Not Discovery)}
\date{2026-11-23}
\begin{document}
\maketitle
\noindent\textbf{Unit convention:} $c=1$, with independent dimensionful $G$ and $\hbar$. All Earth/Moon Yukawa results are illustrative model predictions, not observed deviations. The strong, weak, and electromagnetic forces are not derived here.
\section{A constant-speed helical curve}
Let
\begin{equation}
\mathbf{x}(t)=(a\cos \omega t,a\sin\omega t,ut),\qquad
  a^2\omega^2+u^2=1.
\end{equation}
Its curvature and torsion are
\begin{equation}
\kappa=\frac{a\omega^2}{a^2\omega^2+u^2}=a\omega^2,
\quad \tau=\frac{u\omega}{a^2\omega^2+u^2}=u\omega,
\quad \boxed{\kappa^2+\tau^2=\omega^2}.
\end{equation}
Its nonzero coordinate acceleration is $\ddot{\mathbf x}=(-a\omega^2\cos \omega t,-a\omega^2\sin\omega t,0)$; it is not a free inertial massless trajectory in flat spacetime.
For any ratio $z>0$ choose
\begin{equation}
 u=\frac{z}{\sqrt{1+z^2}},\qquad a=\frac{1}{\omega\sqrt{1+z^2}},
\end{equation}
which yields $\tau/\kappa=z$. Thus this kinematics alone cannot uniquely fix $\alpha_{\mathrm{EM}}$.
\section{A counterexample to universal field orthogonality}
Choose homogeneous parallel electromagnetic fields $\mathbf E=(0,0,4)$ and $\mathbf B=(0,0,3)$ in common Heaviside--Lorentz units with $c=1$. Then $\mathbf E\cdot\mathbf B=12\neq0$. Both $I_1=B^2-E^2$ and $I_2=\mathbf E\cdot\mathbf B$ are Lorentz invariants; a generic ``all three forces must be orthogonal'' rule cannot apply to all electromagnetic configurations.
\section{Finite-size spherical Yukawa source}
For an assumed universal Yukawa kernel,
\begin{equation}
\Phi_Y(r)=-G\alpha_Y\int d^3s\,\rho(s)
  \frac{e^{-|\mathbf r-\mathbf s|/\lambda}}{|\mathbf r-\mathbf s|},
\end{equation}
the spherical exterior $r\ge R$ result is
\begin{align}
\Phi_Y(r)&=-\frac{G\alpha_YQ}{r}e^{-r/\lambda},&
Q&=4\pi\int_0^R ds\,s^2\rho(s)
     \frac{\sinh(s/\lambda)}{s/\lambda}.
\end{align}
With $C=Qe^{-R/\lambda}$ and $M=4\pi\int_0^R\rho(s)s^2ds$,
\begin{equation}
|g(r)|=\frac{GM}{r^2}[1+\epsilon(r)],\qquad
\epsilon(r)=\alpha_Y\frac{C}{M}
      \left(1+\frac{r}{\lambda}\right)e^{-(r-R)/\lambda}.
\end{equation}
\section{Two- and three-altitude discriminants}
Define $A(r)=r^2|g(r)|=GM[1+\epsilon(r)]$. Then
\begin{equation}
\boxed{\mathcal{D}(r_1,r_2)=\frac{A(r_1)}{A(r_2)}-1
=\frac{\epsilon(r_1)-\epsilon(r_2)}{1+\epsilon(r_2)}}.
\end{equation}
This ratio is independent of the overall GM factor, under the spherical approximation.
For $R\le r_0<r_1<r_2$ and $\alpha_Y\ne0$, the additional shape diagnostic is
\begin{equation}
\boxed{\mathcal{Q}_3=\frac{A(r_0)-A(r_1)}{A(r_1)-A(r_2)}
=\frac{f(r_0)-f(r_1)}{f(r_1)-f(r_2)}},\qquad
f(r)=\left(1+\frac r\lambda\right)e^{-(r-r_0)/\lambda}.
\end{equation}
The unknown source amplitude, $G$, $M$, and $\alpha_Y$ cancel. The resulting formula is still model-dependent, requiring $\lambda$, spherical symmetry, and cleaned observations. At $\alpha_Y=0$ the observed difference ratio is undefined.
\section{Radial gradients and nearly circular orbital precession}
\begin{align}
\epsilon'(r)&=-\epsilon(r)\frac{r}{\lambda(r+\lambda)},\\
\Omega^2&=\frac{GM}{r^3}(1+\epsilon(r)),\\
\kappa_r^2&=\frac{GM}{r^3}(1+\epsilon(r)+r\epsilon'(r)),\\
\Delta\varpi&=2\pi\left(\sqrt{\frac{\Omega^2}{\kappa_r^2}}-1\right).
\end{align}
These equations are valid only for stable nearly circular orbits in a spherical Newtonian gravitational field plus the explicitly assumed Yukawa term. They do not include full GR, $J_2$, tides or third bodies.
\section{Status}
No independent fundamental $G$, electron coupling, particle spectrum, gauge group, or particle masses have been derived. No GOCE, GRACE-FO, GRAIL or SLR raw observation has been fitted. All automated tests check formulas and numerical consistency only.
\end{document}
